\section{The Hilbert Nullstellensatz}

The Nullstellensatz is the basic algebraic fact, which we have invoked in the
past to justify various examples, that connects the idea of
the $\spec$ of a ring to classical algebraic geometry.

\subsection{Statement and initial proof of the Nullstellensatz}

There are several ways in which the Nullstellensatz can be stated. Let us
start with the following very concrete version.

\begin{theorem} \label{nullstellensatzoverC}
All maximal ideals in the polynomial ring $R=\mathbb{C}[x_1, \dots, x_n]$ come
from points in $\mathbb{C}^n$. In other words, if $\mathfrak{m} \subset R$ is
maximal, then there exist $a_1, \dots, a_n \in \mathbb{C}$ such that
$\mathfrak{m} = (x_1 - a_1, \dots, x_n - a_n)$.
\end{theorem}


The maximal spectrum of $R=\mathbb{C}[x_1, \dots, x_n]$ is thus identified with
$\mathbb{C}^n$.

We shall now reduce \rref{nullstellensatzoverC} to an easier claim.
Let
$\mathfrak{m}\subset R$ be a maximal ideal. Then there is a map
\[ \mathbb{C} \to R \to R/\mathfrak{m}  \]
where $R/\mathfrak{m}$ is thus a finitely generated $\mathbb{C}$-algebra, as
$R$ is.  The ring $R/\mathfrak{m}$ is also a field by maximality.

We would like to show that $R/\mathfrak{m}$ is a finitely generated
$\mathbb{C}$-vector
space. This would imply that $R/\mathfrak{m}$ is integral over $\mathbb{C}$,
and there are no proper algebraic extensions of $\mathbb{C}$. Thus, if we
prove this, it will follow that the map $\mathbb{C} \to R/\mathfrak{m}$ is an
isomorphism. If $a_i \in \mathbb{C}$ ($1 \leq i \leq n$) is the image of
$x_i $ in $R/\mathfrak{m} =
\mathbb{C}$, it will follow that $(x_1 - a_1, \dots, x_n - a_n) \subset
\mathfrak{m}$, so $(x_1 - a_1, \dots, x_n - a_n)=
\mathfrak{m}$.


Consequently, the Nullstellensatz in this form would follow from the next
claim:

\begin{proposition}
Let $k$ be a field, $L/k$ an extension of fields. Suppose $L$ is a finitely
generated $k$-algebra. Then $L$ is a finite $k$-vector space.
\end{proposition}
This is what we will prove.

We start with an easy proof in the special case:
\begin{lemma}
Assume $k$ is uncountable (e.g. $\mathbb{C}$, the original case of interest).
Then the above proposition is true.
\end{lemma}
\begin{proof}
Since $L$ is a finitely generated $k$-algebra, it suffices to show that $L/k$
is algebraic.
If not, there exists $x \in L$ which isn't algebraic over $k$. So $x$ satisfies
no nontrivial polynomials.
I claim now that the uncountably many elements $\frac{1}{x-\lambda}, \lambda
\in K$ are linearly
independent over $K$. This will be a contradiction as $L$ is a finitely
generated $k$-algebra, hence at most countably dimensional over $k$. (Note that
the polynomial ring is countably dimensional over $k$, and $L$ is a quotient.)

So let's prove this. Suppose not. Then there is a nontrivial linear dependence
\[ \sum \frac{c_i}{x - \lambda_i}  = 0, \quad c_i, \lambda_i \in K. \]
Here the $\lambda_j$ are all distinct to make this nontrivial. Clearing
denominators, we find
\[ \sum_i c_i \prod_{j \neq i } (x- \lambda_j) = 0. \]
Without loss of generality, $c_1 \neq 0$.
This equality was in the field $L$. But $x$ is transcendental over $k$. So we
can think of this as a polynomial ring relation.
Since we can think of this as a relation in the polynomial ring, we see that
doing so, all but the $i =1$ term in the sum is divisible by $x - \lambda_1$
as a polynomial.
It follows that, as polynomials in the indeterminate $x$,
\[ x - \lambda_1 \mid c_1 \prod_{j \neq 1} (x - \lambda_j).  \]
This is a contradiction since all the $\lambda_i$ are distinct.
\end{proof}

This is kind of a strange proof, as it exploits the fact that $\mathbb{C}$ is
uncountable.
This shouldn't be relevant.

\subsection{The normalization lemma}

Let's now give a more algebraic proof.
We shall exploit the following highly useful fact in commutative algebra:

\begin{theorem}[Noether normalization lemma] Let $k$ be a field, and $R =
k[x_1, \dots, x_n]/\mathfrak{p}$ be a finitely generated domain over $k$ (where
$\mathfrak{p}$ is a prime ideal in the polynomial ring).

Then there exists a polynomial subalgebra $k[y_1, \dots, y_m] \subset R$ such
that $R$ is integral over $k[y_1, \dots, y_m]$.
\end{theorem}

Later we will see that $m$ is the \emph{dimension} of $R$.

There is a geometric picture here. Then $\spec R$ is some irreducible algebraic
variety in $k^n$ (plus some additional points), with a smaller dimension than
$n$ if $\mathfrak{p} \neq 0$. Then there exists a \emph{finite map} to $k^m$.
In particular, we can map surjectively $\spec R \to k^m$ which is integral.
The fibers are in fact finite, because integrality implies finite fibers.  (We
have not actually proved this yet.)

How do we actually find such a finite projection? In fact, in characteristic
zero, we just take a
vector space projection $\mathbb{C}^n \to \mathbb{C}^m$. For a ``generic''
projection onto a subspace of the appropriate dimension, the projection will
will do as our finite map. In characteristic $p$, this may not work.

\begin{proof}
First, note that $m$ is uniquely determined as the transcendence degree of the
quotient field of $R$ over $k$.

Among the variables $x_1, \dots, x_n \in R$ (which we think of as in $R$ by an
abuse of notation), choose a maximal subset which is algebraically independent.
This subset has no nontrivial polynomial relations. In particular, the ring
generated by that subset is just the polynomial ring on that subset.
We can permute these variables and assume that
$$\left\{x_1,\dots, x_m\right\}$$ is the maximal subset. In particular, $R$
contains the \emph{polynomial ring} $k[x_1, \dots, x_m]$ and is generated by
the rest of the variables. The rest of the variables are not adjoined freely
though.

The strategy is as follows.  We will implement finitely many changes of
variable so that $R$ becomes integral over $k[x_1, \dots, x_m]$.

The essential case is where $m=n-1$. Let us handle this. So we have
\[ R_0 = k[x_1, \dots, x_m]  \subset R = R_0[x_n]/\mathfrak{p}.  \]
Since $x_n$ is not algebraically independent, there is a nonzero polynomial
$f(x_1, \dots, x_m, x_n) \in \mathfrak{p}$.

We want $f$ to be monic in $x_n$. This will buy us integrality. A priori, this
might not be true. We will modify the coordinate system to arrange that,
though. Choose $N \gg 0$. Define for $1 \leq i \leq m$,
\[ x_i' = x_i + x_n^{N^i}.  \]
Then the equation becomes:
\[ 0 = f(x_1, \dots, x_m, x_n) = f( \left\{x_i' - x_n^{N^i}\right\} , x_n). \]
Now $f(x_1, \dots, x_n, x_{n+1})$ looks like some sum
\[ \sum \lambda_{a_1 \dots b} x_1^{a_1} \dots x_m^{a_m} x_n^{b} , \quad
\lambda_{a_1 \dots b} \in k. \]
But $N$ is really really big. Let us expand this expression in the $x_i'$ and
pay attention to the largest power of $x_n$ we see.
We find that
\[ f(\left\{x_i' - x_n^{N_i}\right\},x_n)
\]
has the largest power of $x_n$ precisely where, in the expression for $f$,
$a_m$ is maximized first, then
$a_{m-1}, $ and so on. The largest exponent would have the form
\[ x_n^{a_m N^m + a_{m-1}N^{m-1} + \dots + b}.	\]
We can't, however, get any exponents of $x_n$ in the expression
\( f(\left\{x_i' - x_n^{N_i}\right\},x_n)\) other than these. If $N$ is super
large, then all these exponents will be different from each other.
In particular, each power of $x_n$ appears precisely once in the expansion of
$f$. We see in particular that $x_n$ is integral over $x_1', \dots, x_n'$.
Thus each $x_i$ is as well.

So we find
\begin{quote}
$R$ is integral over $k[x_1', \dots, x_m']$.
\end{quote}

We have thus proved the normalization lemma in the codimension one case. What
about the general case? We repeat this.
Say we have
\[ k[x_1, \dots, x_m] \subset R.  \]
Let $R'$ be the subring of $R$ generated by $x_{1}, \dots,x_m, x_{m+1}$. The
argument we just gave implies that we can choose $x_1', \dots, x_m'$ such that
$R'$ is integral over $k[x_1', \dots, x_m']$, and the $x_i'$ are
algebraically independent.
We know in fact that $R' = k[x_1', \dots, x_m', x_{m+1}]$.

Let us try repeating the argument while thinking about $x_{m+2}$. Let $R'' =
k[x_1', \dots, x_m', x_{m+2}]$ modulo whatever relations that $x_{m+2}$ has to
satisfy. So this is a subring of $R$. The same argument shows that we can
change variables such that $x_1'', \dots, x_m''$ are algebraically independent
and $R''$ is integral over $k[x_1'', \dots, x_m'']$. We have furthermore that
$k[x_1'', \dots, x_m'', x_{m+2}] = R''$.

Having done this, let us give the argument where $m=n-2$. You will then see
how to do the general case. Then I claim that:
\begin{quote}
$R$ is integral over $k[x_1'', \dots, x_m'']$.
\end{quote}
For this, we need to check that $x_{m+1}, x_{m+2}$ are integral (because these
together with the $x''_i$ generate $R''[x_{m+2}][x_{m+2}]=R$.
But $x_{m+2}$ is integral over this by construction. The integral closure of
$k[x_1'', \dots, x_{m}'']$ in $R$ thus contains
\[ k[x_1'', \dots, x''_m, x_{m+2}] = R''.  \]
However, $R''$ contains the elements $x_1', \dots, x_m'$. But by construction,
$x_{m+1}$ is integral over the $x_1', \dots, x_m'$. The integral closure of
$k[x_1'', \dots, x_{m}'']$ must contain $x_{m+2}$. This completes the proof in
the case $m=n-2$. The general case is similar; we just make several changes of
variables, successively.

\end{proof}
\subsection{Back to the Nullstellensatz}

Consider a finitely generated $k$-algebra $R$ which is a field. We need to
show that $R$ is a
finite $k$-module. This will prove the proposition.
Well, note that $R$ is integral over a polynomial ring $k[x_1, \dots, x_m]$ for
some $m$.
If $m > 0$, then this polynomial ring has more than one prime.
For instance, $(0)$ and $(x_1, \dots, x_m)$. But these must lift to primes in
$R$. Indeed, we have seen that whenever you have an integral extension, the
induced map on spectra is surjective. So
\[ \spec R \to \spec k[x_1, \dots, x_m]  \]
is surjective. If $R$ is a field, this means $\spec k[x_1, \dots, x_m]$ has one
point and $m=0$. So $R$ is integral over $k$, thus algebraic. This implies that
$R$ is finite as it is finitely generated. This proves one version of the
Nullstellensatz.


 Another version of the Nullstellensatz, which is
more precise, says:

\begin{theorem} \label{gennullstellensatz}
Let $I \subset \mathbb{C}[x_1, \dots, x_n]$. Let $V \subset \mathbb{C}^n$ be
the subset of $\mathbb{C}^n$ defined by the ideal $I$ (i.e. the zero locus of
$I$).

Then $\rad(I)$ is precisely the collection of $f$ such that $f|_V = 0$. In
particular,
\[ \rad(I) = \bigcap_{\mathfrak{m} \supset I, \mathfrak{m} \
\mathrm{maximal}} \mathfrak{m}.  \]
\end{theorem}

In particular, there is a bijection between radical ideals and algebraic
subsets of $\mathbb{C}^n$.

The last form of the theorem, which follows from the expression of maximal
ideals in the polynomial ring, is very similar to the result
\[ \rad(I) = \bigcap_{\mathfrak{p} \supset I, \mathfrak{p} \
\mathrm{prime}} \mathfrak{p},  \]
true in any commutative ring. However, this general result is not necessarily
true.

\begin{example}
The intersection of all primes in a DVR is zero, but the intersection of all
maximals is nonzero.
\end{example}

\begin{proof}[Proof of \cref{gennullstellensatz}] 
It now suffices to show that for every $\mathfrak{p} \subset \mathbb{C}[x_1,
\dots, x_n]$ prime, we have
\[ \mathfrak{p} = \bigcap_{\mathfrak{m} \supset I \ \mathrm{maximal}}
\mathfrak{m}  \]
since every radical ideal is an intersection of primes.

Let $R = \mathbb{C}[x_1, \dots,
x_n]/\mathfrak{p}$. This is a domain finitely generated over $\mathbb{C}$. We
want to show that the intersection of maximal ideals in $R$ is zero. This is
equivalent to the above displayed equality.

So fix $f \in R - \left\{0\right\}$. Let $R'$ be the localization $R'= R_f$. Then $R'$ is also an
integral domain, finitely generated over $\mathbb{C}$.	$R'$ has a maximal
ideal $\mathfrak{m}$ (which a priori could be zero). If
we look at the map $R' \to R'/\mathfrak{m}$, we get a map into a field
finitely generated
over $\mathbb{C}$, which is thus $\mathbb{C}$.
The composite map
\[ R \to  R' \to R'/\mathfrak{m}   \]
is just given by an $n$-tuple of complex numbers, i.e. to a point in
$\mathbb{C}^n$ which is even in $V$ as it is a map out of $R$. This corresponds
to a maximal ideal in $R$.
This maximal ideal does not contain $f$ by construction.
\end{proof}

\begin{exercise} 
Prove the following result, known as ``Zariski's lemma'' (which easily implies
the Nullstellensatz): if $k$ is a field, $k'$ a field extension of $k$ which
is a finitely generated $k$-\emph{algebra}, then $k'$ is finite algebraic over
$k$. Use the following argument of McCabe (in \cite{Mc76}):
\begin{enumerate}
\item $k'$ contains a subring $S$ of the form $S= k[ x_1, \dots, x_t]$ where
the $x_1, \dots, x_t$ are algebraically independent over $k$, and $k'$ is
algebraic over the quotient field of $S$ (which is a polynomial ring). 
\item  If $k'$ is not algebraic over $k$, then $S \neq k$ is not a field.
\item Show that there is $y \in S$ such that $k'$ is integral over $S_y$.
Deduce that $S_y$ is a field.
\item Since $\spec (S_y) = \left\{0\right\}$, argue that $y$ lies in every
non-zero prime ideal of $\spec S$. Conclude that $1+y \in k$, and $S$ is a
field---contradiction.
\end{enumerate}
\end{exercise} 

\subsection{A little affine algebraic geometry}

In what follows, let $k$ be algebraically closed, and let $A$ be a finitely generated $k$-algebra.  Recall that $\Specm A$ denotes the set of maximal ideals in $A$.  Consider the natural $k$-algebra structure on $\mathrm{Funct}(\Specm A, k)$.  We have a map
$$A \rightarrow \mathrm{Funct}(\Specm A, k)$$
which comes from the Weak Nullstellensatz as follows.  Maximal ideals $\mathfrak{m}\subset A$ are in bijection with maps $\varphi_\mathfrak{m}:A\rightarrow k$ where $\ker(\varphi_\mathfrak{m})=\mathfrak{m}$, so we define $a\longmapsto [\mathfrak{m}\longmapsto \varphi_\mathfrak{m}(a)]$.  If $A$ is reduced, then this map is injective because if $a\in A$ maps to the zero function, then $a\in \cap\, \mathfrak{m}$ $\rightarrow$ $a$ is nilpotent $\rightarrow$ $a=0$.\\

\begin{definition} A function $f\in \mathrm{Funct}(\Specm A,k)$ is called {\bf algebraic} if it is in the image of $A$ under the above map.  (Alternate words for this are {\bf polynomial} and {\bf regular}.) \end{definition}

Let $A$ and $B$ be finitely generated $k$-algebras and $\phi:A\rightarrow B$ a homomorphism.  This yields a map $\Phi:\Specm B\rightarrow \Specm A$ given by taking pre-images.

\begin{definition} A map $\Phi:\Specm B\rightarrow \Specm A$ is called {\bf algebraic} if it comes from a homomorphism $\phi$ as above.\end{definition}

To demonstrate how these definitions relate to one another we have the following proposition.

\begin{proposition} A map $\Phi:\Specm B\rightarrow \Specm A$ is algebraic if and only if for any algebraic function $f\in \mathrm{Funct}(\Specm A,k)$, the pullback $f\circ \Phi\in \mathrm{Funct}(\Specm B,k)$ is algebraic.\end{proposition}

\begin{proof}
Suppose that $\Phi$ is algebraic.  It suffices to check that the following diagram is commutative:
\[ 
\xymatrix{
\mathrm{Funct}(\Specm A,k) \ar[r]^{-\circ\Phi} & \mathrm{Funct}(\Specm B,k) \\
A \ar[u] \ar[r]_{\phi} & B \ar[u]\\
}
\]
where $\phi:A\rightarrow B$ is the map that gives rise to $\Phi$.

[$\Leftarrow$] Suppose that for all algebraic functions $f\in \mathrm{Funct}(\Specm A,k)$, the pull-back $f\circ\Phi$ is algebraic.  Then we have an induced map, obtained by chasing the diagram counter-clockwise:
$$
\xymatrix{
\mathrm{Funct}(\Specm A,k) \ar[r]^{-\circ\Phi} & \mathrm{Funct}(\Specm B,k) \\
A \ar[u] \ar@{-->}[r]_{\phi} & B \ar[u]\\
}
$$
From $\phi$, we can construct the map $\Phi':\Specm B \rightarrow \Specm A$ given by $\Phi'(\mathfrak{m})=\phi^{-1}(\mathfrak{m})$.  I claim that $\Phi=\Phi'$.  If not, then for some $\mathfrak{m}\in \Specm B$ we have $\Phi(\mathfrak{m})\neq \Phi'(\mathfrak{m})$.  By definition, for all algebraic functions $f\in \mathrm{Funct}(\Specm A,k)$, $f\circ\Phi=f\circ\Phi'$ so to arrive at a contradiction we show the following lemma:\\
Given any two distinct points in $\Specm A=V(I)\subset k^n$, there exists some
algebraic $f$ that separates them.  This is trivial when we realize that any
polynomial function is algebraic, and such polynomials separate points.
\end{proof}

